% ============================================================
%  PhysioMIMO-Net — IEEE Journal of Biomedical and Health
%  Informatics (JBHI) manuscript
%  Authors: Shahzaib Ur Rehman, Dr. Joddat Fatima
%  Target: IEEE JBHI two-column format
% ============================================================
\documentclass[journal]{IEEEtran}

% --- Packages ---
\usepackage{amsmath,amssymb,amsfonts}
\usepackage{algorithmic}
\usepackage{graphicx}
\usepackage{textcomp}
\usepackage{xcolor}
\usepackage{booktabs}
\usepackage{multirow}
\usepackage{url}
\usepackage{cite}
\usepackage{hyperref}
\hypersetup{
  colorlinks=true,
  linkcolor=black,
  citecolor=black,
  urlcolor=blue
}

% ============================================================
\begin{document}

\title{PhysioMIMO-Net: Physics-Informed Multi-Stream Deep
       Learning for Four-Class Intracranial Haemorrhage
       Classification Directly from Raw MIMO Radar Signals}

\author{%
  Shahzaib~Ur~Rehman and Joddat~Fatima%
  \thanks{S. Ur Rehman is with the Department of Computer
    Science, FAST National University of Computer and Emerging
    Sciences (FAST-NUCES), Peshawar, Pakistan.
    (e-mail: p248001@pwr.nu.edu.pk).}%
  \thanks{J. Fatima is with the Department of Computer
    Science, Bahria University H-11 Campus, Islamabad,
    Pakistan.}%
  \thanks{Manuscript received; revised; accepted.}%
}

\markboth{IEEE Journal of Biomedical and Health Informatics}%
         {Rehman and Fatima: PhysioMIMO-Net}

\maketitle

% ============================================================
\begin{abstract}
Intracranial haemorrhage is a neurological emergency in which
outcome is tightly coupled to the speed of diagnosis and the
precision of subtype identification.  Existing microwave-based
classifiers are predominantly binary and rely on Delay-and-Sum
(DAS) beamforming as a preprocessing step, which collapses
the multi-antenna radar tensor into a 2-D power map and
discards the inter-channel timing structure that encodes
haemorrhage depth.  We present PhysioMIMO-Net, the first
end-to-end deep learning architecture to classify all four
clinically relevant intracranial haemorrhage subtypes ---
Healthy, Epidural (EDH), Subdural (SDH), and Intracerebral
(ICH) --- directly from the raw, background-subtracted MIMO
radar tensor without any beamforming or image reconstruction.
The architecture combines four complementary
streams: a multi-scale temporal encoder operating at three
temporal strides, a cross-window encoder whose branch
boundaries are derived analytically from tissue-layer
propagation delays, a frequency-domain encoder on the FFT
magnitude spectrum, and a MIMO Transformer with geometric
positional encoding over the eight-antenna ring.  These
streams are fused and trained jointly using a combination of
cross-entropy and Supervised Contrastive loss.  We train and
evaluate on WaveForge, a new synthetic dataset of 3{,}948
samples generated by a GPU-accelerated FDTD Maxwell solver
with unique per-sample head geometry drawn from published
population distributions.  PhysioMIMO-Net achieves
87.37\% four-class test accuracy (with eight-fold test-time
augmentation), a macro F1 of 0.8733, and a Healthy-class
F1/AUC of 1.000/1.000.  Binary healthy-versus-haemorrhage
performance is 100\% accuracy, F1 1.000, and AUC 1.000.  A
4.64-million-parameter model, it exceeds the only prior
published four-class result (approximately 76\%) by more than
11 percentage points while operating without DAS preprocessing
and on a significantly more geometrically diverse dataset.
These results establish the feasibility of physics-informed,
DAS-free multi-class intracranial haemorrhage classification
from raw MIMO signals.
\end{abstract}

\begin{IEEEkeywords}
intracranial haemorrhage, MIMO radar, microwave brain imaging,
deep learning, physics-informed neural network, transformer,
supervised contrastive learning, FDTD, end-to-end classification
\end{IEEEkeywords}

% ============================================================
\section{Introduction}
\label{sec:intro}

% --- Clinical Motivation ---
Intracranial haemorrhage (ICH) accounts for approximately
10--15\% of all strokes worldwide and carries a 30-day
mortality of 40--50\%, yet its outcome is critically sensitive
to the interval between onset and treatment \cite{Krishnamurthi2013}.
Three anatomically distinct haemorrhage subtypes --- epidural
haematoma (EDH), subdural haematoma (SDH), and intracerebral
haemorrhage (ICH proper) --- differ in their location, time
course, and urgency: EDH is typically arterial and may
deteriorate within minutes; SDH evolves more slowly but
demands surgical evacuation; parenchymal ICH carries the
highest fatality rate and requires immediate neurological
intervention.  All three are indistinguishable at clinical
presentation without neuroimaging, yet access to CT or MRI
scanners is unavailable in ambulances, emergency departments
in low-resource settings, and mass-casualty scenarios.

% --- Microwave Imaging Promise ---
Microwave and ultra-wideband (UWB) radar offer a potential
prehospital imaging modality: wearable antenna arrays can be
fabricated at low cost, operate without ionising radiation,
and are compatible with rapid application by first responders.
The clinical feasibility of microwave-based stroke triage was
established by Persson et al.\ \cite{Persson2014}, who
demonstrated that swept-frequency signals from a wearable
helmet distinguished acute stroke from healthy controls with
83--100\% sensitivity.  Subsequent simulation and deep
learning studies have confirmed that the dielectric contrast
between healthy and pathological brain tissue at 0.5--1.5 GHz
is sufficient to produce discriminable multi-antenna scatter
signatures \cite{Hossain2020,Fhager2019,Alon2021}.

% --- Gap ---
Despite this progress, two critical limitations constrain the
clinical utility of published systems.  First, virtually all
published deep learning classifiers address binary detection
(healthy versus any haemorrhage) rather than subtype
discrimination; only one study has attempted four-class
classification \cite{Yin2021}, and that study employed DAS
image reconstruction followed by a shallow ensemble.  Second,
the dominant preprocessing paradigm --- Delay-and-Sum (DAS)
beamforming --- collapses the raw multi-antenna time-domain
tensor into a 2-D power map that discards inter-channel phase
relationships and differential arrival times.  Because
haemorrhage subtype is determined by depth from skull (EDH:
0--9 mm; SDH: 9--18 mm; ICH: $>$18 mm), the temporal
structure most informative for subtype classification is
precisely what DAS preprocessing degrades.

% --- Contributions ---
This paper makes the following four contributions.

\begin{enumerate}

\item \textbf{PhysioMIMO-Net}: the first end-to-end deep
learning architecture for four-class (Healthy / EDH / SDH /
ICH) intracranial haemorrhage classification from raw MIMO
radar tensors without DAS beamforming or image reconstruction.

\item \textbf{Physics-informed temporal segmentation}: window
boundaries in the cross-window encoder are derived
analytically from two-way electromagnetic propagation delays
through each tissue layer using the same Cole-Cole dielectric
model \cite{Gabriel1996} used in data generation, directly
encoding tissue-layer physics into the network topology.

\item \textbf{WaveForge dataset}: a GPU-accelerated FDTD
Maxwell solver \cite{Yee1966} producing 3{,}948 unique
synthetic samples, each with independently randomised head
geometry drawn from published adult population distributions
\cite{Lynnerup2005}, providing a population-level geometric
diversity absent from prior datasets.

\item \textbf{State-of-the-art four-class performance}: 87.37\%
test accuracy, macro F1 0.8733, and 100\% binary
healthy-versus-haemorrhage accuracy, exceeding all confirmed
prior published results on this task.

\end{enumerate}

The remainder of this paper is organised as follows.
Section~\ref{sec:related} surveys related work.
Section~\ref{sec:dataset} describes WaveForge.
Section~\ref{sec:method} presents the PhysioMIMO-Net
architecture and training procedure.
Section~\ref{sec:experiments} reports experimental results.
Section~\ref{sec:discussion} discusses clinical relevance and
limitations.  Section~\ref{sec:conclusion} concludes.


% ============================================================
\section{Related Work}
\label{sec:related}

\subsection{Clinical Foundations}

The neurophysical basis for microwave brain imaging was
established by Persson et al.\ \cite{Persson2014}, who showed
that a wearable multi-antenna helmet could differentiate acute
stroke from healthy controls in a 45-patient cohort.  Their
study confirmed that the dielectric contrast between healthy
brain parenchyma and pathological tissue is detectable at UWB
frequencies, and motivated the subsequent simulation and
machine learning research programme.

\subsection{DAS-Preprocessing Classifiers}

The dominant computational pipeline in published microwave
brain classifiers applies Delay-and-Sum (DAS) beamforming
before any machine learning step, converting the raw
multi-antenna tensor into a 2-D power map.  Hossain et al.\
\cite{Hossain2020} applied a VGG-16 CNN to DAS-reconstructed
images and achieved $\approx$93\% binary accuracy on a fixed
spherical phantom.  Hasan et al.\ \cite{Hasan2024} extended
this to an explainable wavelet CNN on microwave scalograms,
reporting 95.7\% binary accuracy on a physical phantom and
81.7\% on simulation data; their game-theoretic feature
attribution revealed that discriminative information is
concentrated in the 1.3--1.7 ns time window following
transmission, empirically confirming the physical expectation
that haemorrhage depth governs signal timing.  Critically,
DAS preprocessing discards the inter-channel phase
relationships and differential time-of-arrival structure that
encode depth, precisely the quantities needed for subtype
discrimination.

\subsection{End-to-End Raw-Signal Classifiers}

Fhager et al.\ \cite{Fhager2019} bypassed DAS and applied
subspace-distance learning directly to broadband transmission
coefficients from 1{,}000 varied FDTD head models, finding
that classification AUC approached unity as geometric
diversity of the training set increased.  Alon and
Dehkharghani \cite{Alon2021} fed raw multi-frequency
S-parameter matrices (amplitude and phase) from 666
stochastic 2-D FDTD phantoms directly to deep neural
networks, achieving $>$94\% binary accuracy and AUC of 0.996.
This constitutes the closest conceptual precursor to our
approach: stochastic geometry combined with direct
S-parameter input.  Hedayati et al.\ \cite{Hedayati2024}
demonstrated the feasibility on a real measurement system
using an 8-element Vivaldi array and a 3-D printed phantom,
achieving $>$0.99 sensitivity and 1.65 mm localisation
accuracy with end-to-end classification on raw complex
S-parameters.

\subsection{Dataset Scale and Geometry Variation}

A persistent issue in the field is the use of single fixed
phantom geometries or a small number of discrete variants,
which inflates performance estimates by allowing networks to
exploit invariant skull-reflection timing as an implicit depth
reference.  Fhager et al.\ \cite{Fhager2019} established the
importance of geometry variation for reliable binary detection.
Alon and Dehkharghani \cite{Alon2021} extended this to
stochastic 2-D models.  The WaveForge dataset addresses this
at scale by assigning an independently randomised head
geometry to every sample, drawn from the population
distributions reported by Lynnerup et al.\ \cite{Lynnerup2005}
for adult cranial bone thickness.

\subsection{Multi-Class and Four-Class Gap}

Published deep learning classifiers for microwave brain
imaging have overwhelmingly framed the problem as binary
detection.  Three-class studies (distinguishing healthy,
ischaemic, and haemorrhagic stroke) consolidate all
haemorrhage subtypes under a single label.  To the best of
our knowledge, only one prior study has attempted four-class
classification: Yin et al.\ \cite{Yin2021} applied DAS
preprocessing followed by a shallow MLP+SVM ensemble to
$\approx$400 FDTD samples with three discrete head-size
variants, reporting $\approx$76\% accuracy.  PhysioMIMO-Net
is the first end-to-end architecture to perform four-class
discrimination without DAS, achieving 87.37\% accuracy on a
larger and more geometrically diverse dataset.


% ============================================================
\section{Dataset: WaveForge}
\label{sec:dataset}

\subsection{FDTD Simulation Engine}

WaveForge generates synthetic MIMO radar measurements using a
GPU-accelerated finite-difference time-domain (FDTD) Maxwell
solver \cite{Yee1966} implemented in PyTorch \cite{Paszke2019}.
The solver discretises a $64^3$ voxel domain at 3 mm/cell
(192 mm cube) with a time step of $\Delta t = 5.66$ ps,
satisfying the Courant stability condition, and runs for
$N_t = 700$ time steps.

\subsection{Head Phantom and Haemorrhage Model}

Each simulation uses a unique layered spherical head phantom
with geometry parameters sampled independently from
population-level Gaussian priors.  The head outer radius
$r_\text{head} \sim \mathcal{N}(94, 4^2)$ mm and skull
thickness $t_\text{skull} \sim \mathcal{N}(7, 1.5^2)$ mm
are drawn from the adult cranial bone thickness distribution
reported by Lynnerup et al.\ \cite{Lynnerup2005}.  Tissue
dielectric properties at each frequency are assigned using
the Cole-Cole parametric model from Gabriel et al.\
\cite{Gabriel1996}, the community standard for in-silico
bioelectromagnetics.

Three haemorrhage subtypes are modelled as ellipsoidal
inclusions placed within the appropriate anatomical layer:
EDH between skull and dura mater (depth 0--9 mm from inner
skull); SDH between dura and arachnoid mater (depth 9--18 mm);
and ICH within the brain parenchyma (depth $>18$ mm).
Blood dielectric properties are assigned for three ageing
stages (acute, subacute, chronic), reflecting the progressive
clot coagulation and protein concentration changes that alter
the dielectric constant from $\epsilon_r \approx 61$
(acute fresh blood) to $\epsilon_r \approx 44$ (chronic clot).

\subsection{Antenna Array and Signal Acquisition}

An 8-element ring array is positioned at a radius of 90 mm
around the equatorial plane of the head.  Each element
transmits a UWB Gaussian monocycle pulse centred at 1.0 GHz
with bandwidth 0.5--1.5 GHz in turn while all 8 elements
record simultaneously, yielding a full
$8 \times 8 \times 700$ MIMO time-domain tensor per simulation.
Background subtraction isolates scattered signals:
$\mathbf{s}_\text{scat} = \mathbf{s}_\text{total} -
\mathbf{s}_\text{ref}$,
where $\mathbf{s}_\text{ref}$ is the reference measurement
from a phantom with healthy tissue only and the same geometry.

\subsection{Dataset Statistics}

The WaveForge dataset comprises 3{,}948 samples across four
balanced classes (Healthy, EDH, SDH, ICH), split into
$\approx$3{,}200 training and $\approx$748 test samples.
Each sample is a unique phantom; no geometry is repeated.
The dataset and generation code are available at
\url{https://github.com/shahzaibshazoo/waveforge}
\cite{WaveForge2026}.


% ============================================================
\section{Methodology: PhysioMIMO-Net}
\label{sec:method}

\subsection{Overview}

PhysioMIMO-Net processes the raw background-subtracted MIMO
tensor $\mathbf{X} \in \mathbb{R}^{8 \times 8 \times 700}$
through four parallel encoding streams, fuses the resulting
embeddings, and predicts the haemorrhage class label.  The
architecture contains 4.64M trainable parameters.
Fig.~\ref{fig:arch} illustrates the overall design.

\begin{figure}[!t]
\centering
% Replace with actual figure file before submission
\framebox[0.98\columnwidth]{\parbox{0.9\columnwidth}{%
  \centering\vspace{1.5cm}
  [Architecture diagram: four parallel streams feeding into
   a fusion MLP.  To be inserted as a vector figure.]
  \vspace{1.5cm}}}
\caption{PhysioMIMO-Net architecture.  Four parallel encoding
  streams (Multi-Scale Temporal, Cross-Window, Frequency,
  MIMO Attention) each produce a 256-dimensional embedding
  from the same raw $(8,8,700)$ MIMO tensor.  Embeddings are
  concatenated and projected to the class prediction via a
  two-layer MLP.}
\label{fig:arch}
\end{figure}

\subsection{Stream 1: Multi-Scale Temporal Encoder}

The multi-scale temporal encoder captures discriminative
patterns at three temporal resolutions simultaneously.  The
input tensor is first reshaped from
$\mathbb{R}^{8\times 8 \times 700}$ to
$\mathbb{R}^{64 \times 700}$ by flattening the spatial
antenna dimensions.  Three parallel 1-D CNN branches apply
convolutional filters with strides $s \in \{1, 2, 4\}$ over
the full 700-step time axis:
\begin{equation}
  \mathbf{h}_s = \text{BN}\bigl(\sigma\bigl(
    \text{Conv1D}(\mathbf{X}_{64}; s)\bigr)\bigr),
  \quad s \in \{1, 2, 4\},
  \label{eq:mste}
\end{equation}
where $\text{BN}(\cdot)$ denotes batch normalisation
\cite{Ioffe2015} and $\sigma$ is the ReLU activation.
The fine branch ($s=1$) resolves sub-nanoscale reflections
from skull-dura interfaces, where the two-way delay
$\tau_\text{skull} = 2 t_\text{skull} \sqrt{\epsilon_r^\text{skull}} / c
\approx 0.17$ ns corresponds to $\approx 30$ time steps.
The coarse branch ($s=4$) captures the slower ICH
parenchymal return.  Branch features are concatenated and
projected to a 256-dimensional embedding.

\subsection{Stream 2: Cross-Window Encoder}

The cross-window encoder divides the time axis into three
physics-aligned windows whose boundaries are derived from the
two-way electromagnetic propagation delays through each tissue
layer.  Using the Cole-Cole dielectric values from
Gabriel et al.\ \cite{Gabriel1996} at the centre frequency
$f_0 = 1.0$ GHz, the wave velocity in each layer is
$v_i = c / \sqrt{\epsilon_{r,i}}$, and the two-way delay to
each anatomical boundary is:
\begin{align}
  \tau_{\text{EDH}} &= \frac{2 t_\text{skull}}{v_\text{skull}}
    = \frac{2 t_\text{skull} \sqrt{\epsilon_{r}^\text{skull}}}{c},
  \label{eq:tau_edh}\\
  \tau_{\text{SDH}} &= \tau_{\text{EDH}}
    + \frac{2 d_\text{EDH}}{v_\text{CSF}},
  \label{eq:tau_sdh}\\
  \tau_{\text{ICH}} &= \tau_{\text{SDH}}
    + \frac{2 d_\text{SDH}}{v_\text{brain}}.
  \label{eq:tau_ich}
\end{align}
Converting to discrete time steps at $\Delta t = 5.66$ ps,
the three encoder windows are:
\begin{align*}
  W_1 &: t \in [0, 150] \quad \text{(skull and EDH zone)},\\
  W_2 &: t \in [50, 350] \quad \text{(dura and SDH zone)},\\
  W_3 &: t \in [200, 700] \quad \text{(brain and ICH zone)}.
\end{align*}
Windows overlap deliberately to capture transitions between
layers.  Each window is processed by a dedicated 1-D CNN
branch.  A multi-head self-attention layer with 4 heads then
operates over the three branch outputs, learning the energy
ratios between windows --- a feature critical for
discriminating EDH from SDH, which differ only in which
window contains the dominant scatter return.  The fused
output is projected to 256 dimensions.

\subsection{Stream 3: Frequency-Domain Encoder}

The frequency encoder operates on the magnitude spectrum of
the per-antenna-pair MIMO signal.  For each of the 64
antenna-pair channels, a real FFT is computed over the 700-step
time series, yielding a 351-point magnitude spectrum; bins
$b \in [1, 200]$ are retained (0.25--3.3 GHz at the
simulation sampling rate):
\begin{equation}
  \mathbf{F}_{ij} = \bigl|\mathcal{F}\{\mathbf{x}_{ij}\}
    \bigr|_{1:200},
  \quad i,j \in \{1,\ldots,8\}.
  \label{eq:fft}
\end{equation}
A 2-D CNN processes the resulting $8 \times 8 \times 200$
spectral tensor, capturing the spectral fingerprints of
dura-layer dielectric interference that characterise each
haemorrhage type.  The output is projected to 256 dimensions.

\subsection{Stream 4: MIMO Attention Encoder}

The MIMO attention encoder treats each of the 8 antennas as a
token in a Transformer sequence \cite{Vaswani2017}.  For
antenna $a$, all received signals (8 channels of length 700)
are compressed to a 256-dimensional embedding by a 1-D CNN:
\begin{equation}
  \mathbf{e}_a = f_\text{embed}(\mathbf{X}_{a,:,:}),
  \quad a \in \{1,\ldots,8\}.
  \label{eq:token}
\end{equation}
Geometric positional encoding encodes the angular position of
each antenna on the ring array.  For the $a$-th antenna at
azimuthal angle $\phi_a = 2\pi(a-1)/8$:
\begin{equation}
  \mathbf{p}_a = [\cos\phi_a,\; \sin\phi_a,\;
    \cos 2\phi_a,\; \sin 2\phi_a,\; \ldots].
  \label{eq:posenc}
\end{equation}
This positional encoding reflects the ring geometry and allows
the attention mechanism to learn which antenna pairs provide
complementary phase information for each haemorrhage subtype.
A 4-layer Transformer encoder with 8 attention heads operates
over the sequence of 8 antenna tokens.  The CLS token
embedding (256 dimensions) is used as the stream output.

\subsection{Fusion and Classification Head}

The four 256-dimensional stream embeddings are concatenated
into a 1024-dimensional vector, which is processed by a
two-layer MLP with residual connection:
\begin{equation}
  \mathbf{z} = \text{MLP}_2\bigl(\text{LayerNorm}(
    [\mathbf{h}_\text{MSTE};\,
     \mathbf{h}_\text{CWE};\,
     \mathbf{h}_\text{FE};\,
     \mathbf{h}_\text{MIMO}])\bigr),
  \label{eq:fusion}
\end{equation}
with $\mathbf{z} \in \mathbb{R}^{256}$.  A linear layer maps
$\mathbf{z}$ to the four-class logit vector.

\subsection{Loss Function}

Training minimises a weighted combination of cross-entropy
and Supervised Contrastive Loss \cite{Khosla2020}:
\begin{equation}
  \mathcal{L} = \mathcal{L}_\text{CE}(\hat{y}, y)
    + \lambda_\text{SupCon}\,
    \mathcal{L}_\text{SupCon}(\mathbf{z}, y),
  \label{eq:loss}
\end{equation}
with $\lambda_\text{SupCon} = 0.3$ and temperature
$\tau = 0.3$.  Supervised contrastive loss draws embeddings
of same-class samples together and pushes different-class
embeddings apart in the representation space.  This was found
empirically to equalise the per-class F1 scores for the
hard-to-separate EDH and SDH classes, both converging to
F1 = 0.755 in the final model.

\subsection{Test-Time Augmentation}

At inference time, each test sample undergoes eight
augmentations via cyclic rotation of the antenna ring index
(rotations by 0, 45, 90, \ldots, 315 degrees), exploiting
the ring symmetry of the array.  The eight softmax output
vectors are averaged, reducing variance in the final class
probability estimate \cite{Shanmugam2021}.  TTA improves
four-class accuracy from 86.84\% (standard) to 87.37\%.

\subsection{Training Details}

The model is trained with the Adam optimiser \cite{Kingma2015}
with an initial learning rate of $10^{-4}$ and cosine
annealing schedule.  Batch size is 32.  Training runs for
54 epochs with early stopping on validation macro F1.
All experiments use PyTorch \cite{Paszke2019} on an NVIDIA
T4 GPU (16 GB VRAM) via Kaggle.


% ============================================================
\section{Experiments}
\label{sec:experiments}

\subsection{Experimental Setup}

All experiments use the WaveForge dataset described in
Section~\ref{sec:dataset}, split into 3{,}200 training and
748 test samples with balanced class distribution.  Models are
evaluated on the held-out test set only.  No test data is
used during training or hyperparameter selection.  Reported
accuracy is the top-1 classification accuracy over all
four classes; macro F1 is the unweighted average of
per-class F1 scores; per-class AUC is the area under the
one-versus-rest ROC curve.

\subsection{Four-Class Performance}

Table~\ref{tab:perclass} presents the per-class and aggregate
performance of PhysioMIMO-Net on the four-class test set.

\begin{table}[!t]
\caption{Per-Class and Aggregate Performance of
  PhysioMIMO-Net on the Four-Class Test Set}
\label{tab:perclass}
\centering
\begin{tabular}{lcccc}
\toprule
Class & Precision & Recall & F1 & AUC \\
\midrule
Healthy        & 1.000 & 1.000 & 1.000 & 1.000 \\
Epidural (EDH) & 0.755 & 0.755 & 0.755 & 0.947 \\
Subdural (SDH) & 0.755 & 0.755 & 0.755 & 0.940 \\
Intracerebral (ICH) & 0.984 & 0.984 & 0.984 & 1.000 \\
\midrule
\textbf{Macro avg} & \textbf{0.874} & \textbf{0.874}
  & \textbf{0.874} & \textbf{0.972} \\
\midrule
\multicolumn{5}{l}{Four-class accuracy (standard): 86.84\%} \\
\multicolumn{5}{l}{Four-class accuracy (TTA $\times$8): \textbf{87.37\%}} \\
\multicolumn{5}{l}{Macro F1 (TTA): \textbf{0.8733}} \\
\bottomrule
\end{tabular}
\end{table}

The Healthy and ICH classes are separated with near-perfect
F1 and AUC values.  The EDH and SDH classes achieve identical
F1 = 0.755, with AUC values of 0.947 and 0.940 respectively.
The high AUC values for EDH and SDH indicate strong
discriminative information in the learned representations;
the moderate F1 reflects a tight physical decision boundary
between two classes that differ by only $\approx$9 mm in
depth from skull, at the resolution limit of 0.5--1.5 GHz
UWB signals in high-permittivity tissue.

\subsection{Binary Performance}

The binary task (Healthy versus any haemorrhage) is the
primary clinical screening scenario.  PhysioMIMO-Net achieves
100\% accuracy, F1 = 1.000, and AUC = 1.000 on this binary
test.  This indicates that the model has learned a complete
separation between the Healthy class and all three haemorrhage
subtypes, which is the clinically highest-stakes outcome
(no missed bleeds).

\subsection{Comparison with State of the Art}

Table~\ref{tab:sota} compares PhysioMIMO-Net with published
microwave brain haemorrhage classifiers.  Direct numerical
comparison must be interpreted cautiously because the studies
differ in class count, input modality, preprocessing, and
dataset diversity (Section~\ref{sec:discussion}).

\begin{table*}[!t]
\caption{Comparison with Published Microwave Brain Haemorrhage
  Classification Systems}
\label{tab:sota}
\centering
\begin{tabular}{llllllll}
\toprule
Study & Year & Classes & Input & DAS? & Dataset $N$ &
  Phantom Variety & Best Reported Result \\
\midrule
Persson et al.\ \cite{Persson2014}
  & 2014 & 2 (stroke vs.\ control) & Multi-antenna S-params
  & Yes & 45 (clinical) & Clinical & Sens.\ 83--100\% \\
Fhager et al.\ \cite{Fhager2019}
  & 2019 & 2 & Transmission coefficients
  & No & 1{,}000 FDTD & 1{,}000 unique & AUC $\to$ 1.0 \\
Hossain et al.\ \cite{Hossain2020}
  & 2020 & 2 & DAS images
  & Yes & $\approx$200 FDTD & 1 fixed & $\approx$93.0\% \\
Alon \& Dehkharghani \cite{Alon2021}
  & 2021 & 2 & Raw S-params (amp+phase)
  & No & 666 FDTD & Stochastic 2-D & $>$94\% / AUC 0.996 \\
Yin et al.\ \cite{Yin2021}
  & 2021 & \textbf{4} & DAS images
  & Yes & $\approx$400 FDTD & 3 discrete & $\approx$76.0\% \\
Hasan et al.\ \cite{Hasan2024}
  & 2024 & 2 & Wavelet scalograms
  & Partial & Sim.\ + phantom & Fixed & 95.7\% (phantom) \\
Hedayati et al.\ \cite{Hedayati2024}
  & 2024 & 2 & Raw complex S-params
  & No & Experimental & Robotic phantom & Sens.\ $>$0.99 \\
\midrule
\textbf{PhysioMIMO-Net (ours)}
  & 2026 & \textbf{4} & Raw MIMO tensor
  & \textbf{No} & \textbf{3{,}948 FDTD}
  & \textbf{Unique per sample}
  & \textbf{87.37\% / macro F1 0.8733} \\
\bottomrule
\multicolumn{8}{l}{\footnotesize Yin et al.\ \cite{Yin2021}
  DOI unconfirmed; cited as best prior 4-class result.
  ``DAS'' column reflects the primary input preprocessing.}
\end{tabular}
\end{table*}

PhysioMIMO-Net achieves the highest four-class accuracy by
more than 11 percentage points over the only prior four-class
system, while operating without DAS preprocessing on a dataset
with unique per-sample phantom geometry.

\subsection{Ablation Study}

Table~\ref{tab:ablation} reports an ablation study conducted
on an earlier version of the architecture (v1) trained on a
combined Part~1+Part~2 subset of WaveForge.  The baseline
model achieves 84.71\% four-class accuracy.

\begin{table}[!t]
\caption{Ablation Study: Contribution of Each Architecture
  Component (v1 Architecture, Part~1+Part~2 Dataset)}
\label{tab:ablation}
\centering
\begin{tabular}{lcc}
\toprule
Configuration & Accuracy & $\Delta$ vs.\ Baseline \\
\midrule
Full model (baseline) & 84.71\% & --- \\
Without MIMO Transformer & 79.52\% & $-$5.19 pp \\
Without temporal split (single CNN) & 81.78\% & $-$2.93 pp \\
Without multi-task regression & 83.11\% & $-$1.60 pp \\
\bottomrule
\multicolumn{3}{l}{\footnotesize pp = percentage points.
  All conditions trained identically.}
\end{tabular}
\end{table}

Removing the MIMO Transformer causes the largest accuracy
drop (5.19 pp), confirming that cross-antenna attention is
the single most important inductive bias in the architecture:
the angular diversity of the 8-element ring provides
complementary phase information that a purely channel-stacked
CNN cannot exploit.  Removing the physics-informed temporal
split reduces accuracy by 2.93 pp, validating the hypothesis
that partitioning the time axis according to tissue-layer
propagation delays improves class discrimination.  Removing
the multi-task regression head causes a 1.60 pp reduction,
indicating that depth supervision provides a useful inductive
bias but is less critical than the attention and temporal
splitting components.

\subsection{Per-Class Analysis}

The Healthy class achieves perfect F1 and AUC (1.000/1.000),
indicating complete separation from all haemorrhage classes.
This is the most clinically critical result: a system with
zero false negatives for the Healthy class guarantees no
missed bleeds at the screening stage.

The ICH class achieves F1 = 0.984 and AUC = 1.000.
Intracerebral haemorrhage produces the largest and deepest
scatter volume, with a return arriving in window $W_3$ at
$t > 200$ steps.  The coarse branch of the multi-scale
encoder captures this return distinctively.

The EDH and SDH classes achieve identical F1 = 0.755, with
AUC = 0.947 and 0.940 respectively.  The convergence of F1
scores for these two classes is a direct consequence of
Supervised Contrastive training: prior to SupCon, EDH
achieved higher F1 than SDH, reflecting a bias towards the
shallower, higher-contrast EDH scatter.  SupCon training
equalised the per-class penalty and produced symmetric
precision and recall for both EDH and SDH.


% ============================================================
\section{Discussion}
\label{sec:discussion}

\subsection{EDH/SDH Equalisation by Supervised Contrastive
            Training}

The most physically challenging aspect of the four-class
problem is discriminating EDH from SDH.  These two subtypes
differ only in whether the haemorrhage is between skull and
dura mater (EDH, depth $\approx$0--9 mm) or between dura
and arachnoid mater (SDH, depth $\approx$9--18 mm).  At
0.5--1.5 GHz, the free-space wavelength at 1 GHz is 300 mm;
in brain tissue with $\epsilon_r \approx 40$, the wavelength
is approximately 47 mm.  A 9 mm depth separation is thus
considerably smaller than one wavelength, approaching the
fundamental resolution limit of this imaging modality.

Without Supervised Contrastive loss, the classifier exhibited
a bias towards EDH, which produces a higher-contrast scatter
due to the abrupt skull-air boundary.  SupCon training
enforces representation space compactness within each class
while maximising separation between classes at the embedding
level.  This causes the model to find features that are
equally informative for both EDH and SDH, resulting in the
symmetric F1 = 0.755 observed in both classes.  The high AUC
values (0.947 and 0.940) confirm that the underlying
discriminative information exists in the raw signals; the
moderate F1 reflects a physically tight decision boundary
rather than an undertrained model.

\subsection{Clinical Relevance}

The binary result (100\% accuracy, F1 1.000, AUC 1.000)
indicates that PhysioMIMO-Net can perfectly separate healthy
brains from any haemorrhage in this synthetic test set.  In
a clinical screening context, this corresponds to a system
that refers all haemorrhage cases for further evaluation while
clearing healthy patients, which is the most safety-critical
operating regime.

The four-class result provides the finer discrimination needed
for triage and treatment decisions: EDH and SDH may both
require surgical evacuation, but EDH is typically more urgent
due to its arterial origin; ICH management differs
substantially from subdural or epidural management.  An
87.37\% four-class system would, in principle, correctly
classify approximately 88 in 100 patients, with the
misclassifications concentrated at the physically difficult
EDH/SDH boundary rather than at the clinically most dangerous
Healthy/Haemorrhage boundary.

\subsection{Limitations and Future Work}

This work constitutes a simulation feasibility study
establishing that end-to-end four-class classification from
raw MIMO signals is achievable in principle.  The following
limitations must be acknowledged before clinical translation.

\textbf{Synthetic data only.}  All results are obtained on
FDTD-simulated data.  The Cole-Cole tissue dielectric model
\cite{Gabriel1996} is the community standard and has been
validated against measurements across the relevant frequency
range, but simulation-to-measurement transfer entails
modelling gaps: antenna-to-skin coupling, body motion
artefacts, inter-individual tissue property variation
(particularly at the skull-brain interface), and
electromagnetic coupling artefacts from metallic implants.
Physical phantom validation and clinical transfer are
identified as the immediate next development milestones,
following the precedent established in this domain by
Fhager et al.\ \cite{Fhager2019} and Hedayati et al.\
\cite{Hedayati2024}.

\textbf{No DAS baseline on the same dataset.}  The claim that
bypassing DAS preserves information critical for subtype
classification is physically motivated and supported by the
ablation study, but a direct numerical comparison between
PhysioMIMO-Net and a DAS-CNN trained on the same WaveForge
data has not been conducted.  This comparison is planned for
future work.

\textbf{Ischaemic stroke not modelled.}  The WaveForge dataset
models haemorrhagic subtypes only.  A clinically complete
system must also discriminate ischaemic stroke, which
produces a different and smaller dielectric contrast.
Extending the class set to five categories
(Healthy / EDH / SDH / ICH / Ischaemic) is a natural next
step.

\textbf{Simplified head geometry.}  The layered spherical
phantom is a first-order approximation of the human head.
Incorporating cortical gyration, CSF spaces, and lateral
brain structure from MRI-derived segmented models is expected
to increase simulation realism at the cost of computational
load.


% ============================================================
\section{Conclusion}
\label{sec:conclusion}

We have presented PhysioMIMO-Net, the first end-to-end deep
learning system for four-class intracranial haemorrhage
classification (Healthy / EDH / SDH / ICH) directly from raw
MIMO radar tensors without Delay-and-Sum beamforming or image
reconstruction.  By combining physics-informed temporal
segmentation, multi-scale temporal encoding, frequency-domain
spectral analysis, and a geometric MIMO Transformer, the
architecture exploits the full information content of the
background-subtracted $(8,8,700)$ MIMO tensor in a
principled, physically motivated manner.  Trained with
Supervised Contrastive Loss on the WaveForge FDTD dataset ---
3{,}948 samples each with a unique randomised head geometry ---
the model achieves 87.37\% four-class test accuracy
(macro F1 0.8733) and 100\% binary healthy-versus-haemorrhage
accuracy, exceeding all confirmed prior published results.
The dominant architectural contribution is the MIMO
Transformer (ablation: $-$5.19 pp without it), confirming
the value of cross-antenna attention for exploiting geometric
diversity in the ring array.  Future work will validate the
system on physical phantoms and extend the class set to
include ischaemic stroke.


% ============================================================
\bibliographystyle{IEEEtran}
\bibliography{refs}

\end{document}
